% ===================== 第 1 章 考试概况 =====================
\section{考试概况}

\subsection{题型结构与分值}

CET-6 总分 710 分，不设及格线，但通常以 \textbf{425 分}作为“通过”的参考标准；四级 425 分以上方可报考六级。笔试时长共 130 分钟，题型与分值如下：

\begin{center}
\small
\begin{longtable}{p{2.6cm}p{3.4cm}p{1.6cm}p{1.9cm}p{4.6cm}}
\toprule
\textbf{部分} & \textbf{题型} & \textbf{题量} & \textbf{分值} & \textbf{时间 / 说明} \\
\midrule
\endfirsthead
\toprule
\textbf{部分} & \textbf{题型} & \textbf{题量} & \textbf{分值} & \textbf{时间 / 说明} \\
\midrule
\endhead
写作 & 短文写作 & 1 题 & 106.5（15\%） & 30 分钟，150--200 词 \\
\midrule
听力 & 长对话 & 8 题 & 56.8（8\%） & 每题 7.1 分 \\
 & 听力篇章 & 7 题 & 49.7（7\%） & 每题 7.1 分 \\
 & 讲座 / 讲话 & 10 题 & 142（20\%） & 每题 14.2 分，占听力近六成 \\
\midrule
阅读 & 选词填空 & 10 题 & 35.5（5\%） & 每题 3.55 分，15 选 10 \\
 & 长篇阅读（匹配） & 10 题 & 71（10\%） & 每题 7.1 分 \\
 & 仔细阅读 & 10 题 & 142（20\%） & 每题 14.2 分，2 篇 \\
\midrule
翻译 & 汉译英段落 & 1 题 & 106.5（15\%） & 与阅读共享 70 分钟，180--200 汉字 \\
\bottomrule
\end{longtable}
\end{center}

\begin{tipbox}[title=性价比排序（按“每分钟得分”估算）]
仔细阅读（14.2 分/题）> 讲座讲话（14.2 分/题）> 长篇阅读（7.1 分/题）> 写作 $\approx$ 翻译（106.5 分/部分）> 选词填空（3.55 分/题）。复习时间紧张时，优先抓\textbf{仔细阅读、翻译和写作}这三块。
\end{tipbox}

\subsection{考场流程（以官方通知为准）}

典型的下午考试流程如下，实际时间以考点安排为准：

\begin{center}
\small
\begin{longtable}{p{3.2cm}p{10.2cm}}
\toprule
\textbf{时间} & \textbf{环节} \\
\midrule
\endfirsthead
\toprule
\textbf{时间} & \textbf{环节} \\
\midrule
\endhead
14:40--15:00 & 考生入场、试音（耳机调频至考点频率） \\
15:00--15:10 & 发答题卡与试卷，填写个人信息，禁止提前答题 \\
15:10--15:40 & 写作（答题卡 1） \\
15:40--16:10 & 听力（答题卡 1），\textbf{听力结束立即收答题卡 1} \\
16:15--17:25 & 阅读 + 翻译（答题卡 2） \\
17:25 & 考试结束，收答题卡 2 与试题册 \\
\bottomrule
\end{longtable}
\end{center}

\begin{notebox}[title=三条“保命”规则]
\begin{enumerate}
  \item \textbf{听力必须边听边涂卡}。听力播放结束后立刻收答题卡 1，没有单独涂卡时间。
  \item 写作和听力共用答题卡 1；阅读和翻译共用答题卡 2 且共享 70 分钟，务必给翻译留出至少 25 分钟。
  \item 答题顺序可以自定：推荐\textbf{仔细阅读 $\to$ 翻译 $\to$ 长篇阅读 $\to$ 选词填空}，或按个人习惯先写翻译再读文章，避免最后时刻手忙脚乱。
\end{enumerate}
\end{notebox}

\subsection{近年命题趋势（写作与翻译）}

\begin{itemize}
  \item \textbf{写作}：以引言式议论文为主（给出首句或名言，续写成文），话题集中在个人成长、教育、科技与人际关系；“现象解释型”“问题解决型”“观点选择型”“谚语评述型”四种框架轮换出现。
  \item \textbf{翻译}：全部为“中国主题”段落（180--200 汉字），文化类约占四成，社会经济类约三成半，地理生活类约两成半；近年趋向“传统文化 + 时代热点”的捆绑考查（如非遗 + 短视频、传统技艺 + 数字化保护）。
  \item \textbf{选词填空}：15 个选项中名词 3--5 个、动词 5--7 个（含三单 / -ed / -ing 形式）、形容词 2--5 个、副词 2--3 个，选项词均为四类实词，且\textbf{每个词只能使用一次}。
  \item \textbf{听力}：讲座 / 讲话（Section C）分值最重，背景多为社科、科普类话题；长对话以校园、职场、生活场景为主。
\end{itemize}
